\documentclass[a4paper]{article}
% Español (México) con XeLaTeX: fontspec para fuentes Unicode, polyglossia
% para la división silábica y los nombres del documento en la variante mexicana.
\usepackage{fontspec}
\usepackage{polyglossia}
\setdefaultlanguage[variant=mexican]{spanish}
% Los nombres chinos van como «pinyin (original)»: xeCJK da a los caracteres
% Han su propia fuente; sin ella XeLaTeX los omite con un aviso.
% FandolSong no cubre algunos caracteres de nombres (U+73FA); WenQuanYi Zen Hei sí.
\usepackage[AutoFallBack=true]{xeCJK}
\setCJKmainfont{FandolSong-Regular.otf}
\setCJKfallbackfamilyfont{\CJKrmdefault}{WenQuanYi Zen Hei}
% polyglossia no traduce los nombres de listings.
\AtBeginDocument{\providecommand{\lstlistingname}{}\renewcommand{\lstlistingname}{Listado}%
  \providecommand{\lstlistlistingname}{}\renewcommand{\lstlistlistingname}{Índice de listados}}
\usepackage{amsmath, amssymb}
\usepackage{graphicx}
\usepackage[margin=2.25cm]{geometry}
\usepackage[most]{tcolorbox}
\usepackage{booktabs}
\usepackage{tabularx}
\usepackage{array}
\usepackage{float}
\usepackage{xcolor}
\usepackage{hyperref}
\usepackage{fancyhdr}
\setCJKmainfont[AutoFakeBold=2.4]{AR PL UMing CN}
\setCJKsansfont[AutoFakeBold=2.4]{Droid Sans Fallback}
\setCJKmonofont[AutoFakeBold=2.4]{Droid Sans Fallback}
\hypersetup{colorlinks=true, linkcolor=blue!55!black, urlcolor=blue!60!black}
\graphicspath{{./}{figures/}}
\setlength{\parskip}{0.55em}
\setlength{\parindent}{2em}
\setlength{\emergencystretch}{3em}
\renewcommand{\arraystretch}{1.18}
\newtcolorbox{knowledgebox}[1]{enhanced, colback=blue!5!white, colframe=blue!70!black, colbacktitle=blue!70!black, coltitle=white, fonttitle=\bfseries, title={#1}, sharp corners, breakable}
\newtcolorbox{importantbox}[1]{enhanced, colback=yellow!10!white, colframe=yellow!75!black, colbacktitle=yellow!75!black, coltitle=black, fonttitle=\bfseries, title={#1}, sharp corners, breakable}
\newtcolorbox{warningbox}[1]{enhanced, colback=red!5!white, colframe=red!70!black, colbacktitle=red!70!black, coltitle=white, fonttitle=\bfseries, title={#1}, sharp corners, breakable}
\pagestyle{fancy}
\fancyhf{}
\fancyfoot[C]{\thepage}
\fancyfoot[R]{\footnotesize\color{gray}\href{https://github.com/hqhq1025/ai-course-notes}{AI Course Notes}}
\renewcommand{\headrulewidth}{0pt}
\renewcommand{\footrulewidth}{0pt}
\newcommand{\videourl}{https://www.youtube.com/watch?v=VdWEE6vOYRw}
\newcommand{\repourl}{https://github.com/hqhq1025/ai-course-notes}

\begin{document}
\begin{titlepage}
\centering
{\Large Notas didácticas de una entrevista a fondo\par}
\vspace{1.0cm}
{\huge\bfseries Xiao Hong (肖弘), fundador de Manus: emprender con aplicaciones de IA, el primer disparo de los Agent y el juego estratégico del Founder}\par
\vspace{0.5cm}
{\Large Zhang Xiaojun Podcast conversa con Xiao Hong\par}
\vspace{0.35cm}
{\large Elaboradas a partir del video público, la estructura de capítulos, la transcripción, la descripción del video y diagramas conceptuales propios\par}
\vspace{0.3cm}
{\large 2025-03-02; incluye dos entrevistas sucesivas, de otoño de 2024 y de primavera de 2025\par}
\vspace{0.85cm}
\includegraphics[width=0.88\textwidth,height=0.42\textheight,keepaspectratio]{cover.jpg}\par
\vfill
\begin{tcolorbox}[width=0.92\textwidth, colback=black!2!white, colframe=black!55, sharp corners]
\textbf{Título del video}: 95. Entrevista de 3 horas con Xiao Hong, fundador de Manus: el mundo no es una extrapolación lineal; sé una variable importante en el juego\par
\textbf{Canal del video}: Zhang Xiaojun Podcast\par
\textbf{Duración del video}: 03:22:47\par
\textbf{Enlace del video}: \href{\videourl}{\nolinkurl{\videourl}}\par
\textbf{Nota sobre los materiales}: YouTube no ofrece subtítulos descargables; el texto se elaboró a partir de una transcripción por bloques con faster-whisper large-v3 en CPU con int8, la descripción del video, la estructura de capítulos y diagramas didácticos propios. Las tomas fijas de la entrevista no se reutilizan en el cuerpo del texto.\par
\textbf{Página del proyecto}: \href{\repourl}{\nolinkurl{\repourl}}\par
\end{tcolorbox}
\end{titlepage}

\begingroup
\small
\setlength{\parskip}{0pt}
\renewcommand{\baselinestretch}{0.92}\selectfont
\tableofcontents
\endgroup
\newpage
\section{Introducción: por qué este episodio es una nota de ruta para emprender con aplicaciones de IA}

Esta sección sitúa primero el episodio. EP95 es una «entrevista en relevo»: la primera parte se grabó en el otoño de 2024, cuando Xiao Hong (肖弘) acababa de cerrar una ronda de financiamiento, y trata sobre el emprendimiento en serie, el capital, Monica y la metodología de las aplicaciones de IA; la segunda parte se grabó después del Año Nuevo chino de 2025, cuando DeepSeek había transformado el ecosistema de base de las aplicaciones de IA en China, y Xiao Hong empieza a hablar de Manus, que estaba por lanzarse y que fue el «primer disparo» de los productos Agent en China. Por ello, este episodio no es una historia estática de un fundador, sino un registro continuo de las reflexiones de un emprendedor de aplicaciones de IA en medio de un cambio acelerado en las capacidades de los modelos.

\begin{importantbox}{Tesis central del episodio}
La línea principal de Xiao Hong puede resumirse en tres frases: primero, una empresa de aplicaciones de IA debe anticipar las capacidades de los modelos y los movimientos de las grandes empresas, y preparar con antelación el contenedor de producto; segundo, las capacidades de los modelos tenderán a la igualación tecnológica, pero la barrera de entrada de una empresa de aplicaciones proviene del escenario de uso, la marca, la comprensión de los usuarios, la retroalimentación y la ejecución; tercero, el mundo no es una extrapolación lineal: el Founder debe pensar en términos de teoría de juegos y convertirse en una variable relevante de la situación.
\end{importantbox}

\begin{knowledgebox}{Nota sobre la estrategia visual}
Este video es una entrevista con plano fijo, sin proyección de pantalla ni demostración de producto. El cuerpo de la nota usa la portada solo para identificar la fuente; todas las imágenes del cuerpo son diagramas conceptuales propios que explican la estructura de la entrevista en relevo, el costo del VC, la anticipación de los movimientos de las grandes empresas, Monica, la clasificación de las aplicaciones de IA, el producto Agent de Manus y el pensamiento estratégico del Founder.
\end{knowledgebox}

\begin{figure}[H]
\centering
\includegraphics[width=0.96\textwidth]{xiaohong-relay-interview.png}
\caption{Mapa de la entrevista en relevo con Xiao Hong: las dos entrevistas registran las reflexiones continuas de un emprendedor de aplicaciones de IA en un entorno inestable. Diagrama conceptual propio, elaborado a partir de la estructura de capítulos del video completo.}
\end{figure}

\begin{knowledgebox}{Lectura de la figura: no es una línea recta, sino dos cortes de estado}
La figura va desde después del financiamiento del otoño de 2024, pasando por el primer emprendimiento, Monica y el impacto de DeepSeek en el Año Nuevo chino de 2025, hasta Manus. Al leerla, conviene notar que un mismo Founder enfrenta, en distintos momentos, capacidades externas de los modelos, entornos de capital, oportunidades de producto y problemas de organización diferentes; por eso sus juicios también cambian.
\end{knowledgebox}

\subsection{Ruta de lectura}

Esta sección presenta la ruta de lectura. El episodio completo puede dividirse en seis preguntas: primera, ¿qué lecciones de producto, de monetización y de capital le dejó a Xiao Hong su primer emprendimiento? Segunda, ¿por qué «el VC es un medio muy caro de obtener financiamiento»? Tercera, ¿por qué Monica eligió la extensión de navegador, el mercado internacional y el all-in-one AI assistant? Cuarta, ¿cómo se clasifican las aplicaciones de IA: complemento del escenario principal, cambio en las capacidades de los modelos y desbordamiento desde dominios verticales? Quinta, ¿por qué se entiende a Manus como el primer disparo de los Agent en China? Sexta, ¿por qué el Founder debe pasar del razonamiento lógico al pensamiento de teoría de juegos?

\noindent\begin{tabularx}{\textwidth}{p{0.20\textwidth}p{0.34\textwidth}X}
\toprule
Pregunta de lectura & Material de la entrevista & Juicio que se debe formar \\
\midrule
Qué aprendió del primer emprendimiento & Herramientas para cuentas oficiales de WeChat, SCRM, monetización, venta de la empresa & Un producto de tipo herramienta debe monetizarse pronto, sin dejarse desviar por grandes narrativas. \\
Cómo influye el capital en las decisiones & El VC es caro, el peso de los inversionistas, la dimensión del capital & El financiamiento es una herramienta; no es el jefe ni es la estrategia en sí. \\
Por qué se creó Monica & Extensión, mercado internacional, ChatGPT for Google, retroalimentación de usuarios & Una aplicación de IA debe ocupar el contenedor de producto que surge cuando las capacidades del modelo se desbordan. \\
Cómo encuentra oportunidades una empresa de aplicaciones & Jasper, ChatGPT, Cursor, DeepSeek, Manus & Observar la brecha entre las capacidades del modelo y el escenario principal del usuario. \\
El significado de Manus como producto & Tareas, herramientas, entorno, memoria, entrega de resultados & El Agent pasa de la conversación a un sistema capaz de hacer cosas. \\
Cómo piensa el Founder & No linealidad, teoría de juegos, la edad frente a la época, codicia, ira e ignorancia & Emprender no es un razonamiento lineal; hay que cambiar la situación. \\
\bottomrule
\end{tabularx}

\begin{warningbox}{Límites temporales}
La primera entrevista se realizó en el otoño de 2024. En esa parte, «este año» se refiere casi siempre a 2024 y «el año pasado», a 2023. La segunda entrevista se realizó después del Año Nuevo chino de 2025 y trata la situación posterior al auge de DeepSeek y previa al lanzamiento de Manus.
\end{warningbox}

\subsection{Resumen de la sección}

EP95 es una nota de metodología para emprender con aplicaciones de IA. Puede leerse junto con la última entrevista de EP128 sobre Peak/Manus, EP101 sobre YouWare y EP110 sobre el informe técnico de Agent: aquella describe la organización de Manus en una etapa posterior y su situación antes de la adquisición, mientras que este episodio recoge las reflexiones del fundador antes del lanzamiento de Manus y la metodología de base para emprender con aplicaciones.
\section{Primera empresa: de las herramientas para WeChat a las lecciones de comercialización}

La sección anterior trazó la ruta; esta sección examina primero la primera empresa de Xiao Hong (肖弘). En 2015, al egresar de la Universidad de Ciencia y Tecnología de Huazhong, emprendió de inmediato. Al principio desarrolló productos para el entorno universitario que tuvieron mala respuesta; después pasó a un editor para cuentas oficiales de WeChat y, más tarde, a un SCRM para WeCom (企业微信). La formación más importante de esta etapa no fue solo haber creado productos, sino aprender del fracaso, del flujo de efectivo, de la comercialización, del capital y de la venta de la empresa: en un producto de herramienta, una historia más grande no es mejor; lo que hace falta es encontrar usuarios reales, pagos reales y una economía unitaria sostenible.

\begin{figure}[H]
\centering
\includegraphics[width=0.92\textwidth]{founder-cycle.png}
\caption{Ciclo del emprendedor serial: en habilidades, capital, mercado, producto y mentalidad aparecen continuamente nuevas dimensiones. Diagrama conceptual propio, elaborado a partir de la conversación entre 00:05:15 y 00:58:15.}
\end{figure}

\begin{knowledgebox}{Lectura de la figura: la capacidad emprendedora no se completa de una sola vez}
Alrededor del Founder, la figura muestra los puntos de habilidad, el capital, el mercado, el producto, la mentalidad y la negociación estratégica. Xiao Hong repite que cada vez que sentía que sus puntos de habilidad ya eran suficientes, surgía una dimensión nueva. En la primera empresa aprendió primero de producto; después, de comercialización, y más adelante descubrió que el capital también era una dimensión nueva.
\end{knowledgebox}

\subsection{De la retroalimentación positiva en la universidad a la negativa en el mundo real}

Esta subsección explica los primeros reveses. En la universidad, Xiao Hong participó en un club técnico, en un centro de nuevos medios, en un blog y en proyectos de producto, y recibió mucha retroalimentación positiva. Al egresar, desarrolló productos como una red social anónima y un mercado de artículos de segunda mano, pero no consiguieron usuarios, y el equipo estuvo a punto de irse a trabajar a Pekín. Esta experiencia muestra que los pequeños éxitos en el entorno universitario no pueden extrapolarse directamente al mercado real. Emprender de verdad es más duro, porque los usuarios no son el círculo de conocidos, los canales no existen de forma natural y el pago no ocurre por sí solo.

\begin{warningbox}{Error común: tomar la retroalimentación positiva de la universidad como validación de mercado}
Crear herramientas, cuentas oficiales o pequeños productos en la universidad demuestra capacidad e interés personales, pero no demuestra que exista un mercado comercial. La primera lección después de emprender suele ser descubrir que lo que se había recibido era retroalimentación positiva local, y no una demanda amplia.
\end{warningbox}

\subsection{Por qué la comercialización debe empezar antes}

Esta subsección examina la comercialización del editor para cuentas oficiales. En su retrospectiva, Xiao Hong dice que el equipo dedicó demasiado tiempo al crecimiento de usuarios y a la iteración del producto, y pensó poco en la comercialización. Más tarde decidieron vender software, como un SaaS, en lugar de crear una plataforma de intermediación publicitaria. La intermediación publicitaria parece tener un techo más alto, pero no necesariamente forma parte de las capacidades centrales de un equipo de software; vender software parece más pequeño, pero es más fiel a lo que el equipo sabe hacer, cierra mejor el ciclo y además permite al equipo mejorar el producto de forma continua.

\begin{importantbox}{Experiencia práctica: un techo alto no significa que se pueda lograr}
Muchos equipos emprendedores, para contar una historia más grande, eligen modelos de negocio que no corresponden a sus capacidades. En el caso del editor para cuentas oficiales, «captar clientes con la herramienta y después hacer intermediación publicitaria» sonaba más grande, pero lo que realmente funcionó fue el software por suscripción. En un producto de herramienta, que los usuarios paguen directamente por el software forma, de hecho, un ciclo cerrado más sano.
\end{importantbox}

\subsection{El VC es un medio costoso de obtener fondos}

Esta subsección entra en el tema del capital. Xiao Hong dice que lo costoso del VC no se nota cuando a la empresa le va mal, sino cuando le va bien: si la empresa tiene un gran éxito, el costo de la financiación temprana con capital accionario se multiplica. Financiarse es una herramienta razonable, pero el fundador debe saber que cambiará la historia, los objetivos y la presión. Recaudar demasiado dinero puede obligar a la empresa a contar una historia más grande y a hacer cosas que no necesariamente son correctas.

\begin{figure}[H]
\centering
\includegraphics[width=0.90\textwidth]{vc-expensive-capital.png}
\caption{El VC es un medio costoso de obtener fondos: lo costoso no aparece en los momentos difíciles, sino cuando a la empresa le va bien. Diagrama conceptual propio, elaborado a partir de la conversación entre 00:33:16 y 00:42:10.}
\end{figure}

\begin{knowledgebox}{Lectura de la figura: el capital es una herramienta, no una dirección}
A la izquierda, en los momentos difíciles, la financiación parece dinero de rescate; a la derecha, cuando a la empresa le va bien, la dilución temprana y los compromisos mayores se multiplican. Al leer esta figura conviene retener lo siguiente: el VC no es malo; lo esencial es que el Founder sepa qué restricciones obtiene a cambio de usar esta herramienta.
\end{knowledgebox}

\subsection{El inversionista no es el jefe}

Esta subsección conserva una founder voice de mucho valor. Xiao Hong recomienda no tratar a los inversionistas como jefes: si un inversionista comparte un artículo, ese gesto debería pesar más o menos como la opinión de un KOL, y no convertirse automáticamente en estrategia de la empresa. Si un inversionista de verdad tiene una opinión firme, debe convocarse una reunión formal y dejar registro de quién sostiene qué juicio, en lugar de dejar que el Founder adivine si el inversionista quiere que la empresa cambie de rumbo. La razón es que el apalancamiento del Founder es altísimo: un solo error de juicio arrastra el dinero y la dirección del equipo.

\begin{importantbox}{Nota de clase: las decisiones de alto apalancamiento requieren un mecanismo formal}
El Founder puede escuchar la opinión de los inversionistas, pero no debe tomar una sugerencia casual como una orden. Las decisiones importantes de rumbo deben pasar a una discusión formal, con opiniones, evidencia, responsabilidades y mecanismo de decisión explícitos. De lo contrario, el mayor peligro no es que el inversionista se equivoque, sino que el Founder malinterprete un comentario dicho de pasada.
\end{importantbox}

\subsection{Resumen de la sección}

La primera empresa le dejó a Xiao Hong tres lecciones de fondo: en un producto de herramienta, la comercialización debe empezar pronto; la financiación es una herramienta costosa, no la estrategia en sí, y el Founder debe conservar la soberanía de su juicio. Estas lecciones influyeron después directamente en sus decisiones sobre Monica, la expansión internacional y Manus.
\section{Anticipar lo que anticiparán las grandes empresas: espera, tentación y punto de inflexión}

La sección anterior trató el capital y la comercialización; esta examina la capacidad táctica más importante de Xiao Hong (肖弘): anticipar lo que anticiparán las grandes empresas. La oportunidad del SCRM para WeCom (企业微信) surgió de un juicio: Tencent (腾讯) iba a regular los complementos no oficiales de WeChat personal y, al mismo tiempo, iba a canalizar la demanda de cumplimiento normativo hacia el ecosistema de WeCom. No se sabía exactamente cuándo ocurriría, pero era posible prepararse con anticipación dentro de un costo asumible. Cuando se combatieron los complementos no oficiales, el producto ya preparado pudo atender de inmediato esa demanda.

\begin{figure}[H]
\centering
\includegraphics[width=0.96\textwidth]{predict-giants.png}
\caption{Anticipar lo que anticiparán las grandes empresas: la espera, la tentación y el anticonsenso al estilo de «La gran apuesta» forman en conjunto el punto de inflexión decisivo. Diagrama conceptual propio, elaborado a partir de la conversación entre 00:42:10 y 00:52:27.}
\end{figure}

\begin{knowledgebox}{Lectura de la figura: se sabe que ocurrirá, pero no cuándo}
La figura avanza de la observación a la anticipación, la espera, la tentación y el punto de inflexión. El método central de Xiao Hong es este: si juzgas que algo ocurrirá con alta probabilidad, pero no conoces el momento, prepárate con anticipación dentro de un costo asumible. La dificultad del emprendimiento está en que durante la espera no hay retroalimentación positiva, y además aparecen proyectos a la medida, pedidos de corto plazo y titubeos en el equipo.
\end{knowledgebox}

\subsection{La espera al estilo de «La gran apuesta»: por qué es lo más difícil}

Esta subsección explica lo doloroso de la espera. Después del lanzamiento del producto SCRM, los complementos no oficiales todavía no se habían regulado, el producto no tenía usuarios y el equipo tendía a dudar de la estrategia. Al mismo tiempo surgían tentaciones: desarrollos a la medida para grandes clientes, ingresos de corto plazo y otras direcciones. Xiao Hong recurre a la analogía de «La gran apuesta»: acertar en el juicio no equivale a ganar dinero de inmediato; durante la espera hay que soportar presión de efectivo, de equipo y de confianza. Finalmente, cuando se bloquearon los complementos no oficiales, esa misma noche escribieron un artículo, lo difundieron con publicidad pagada y se posicionaron rápidamente en la mente de los usuarios como la alternativa.

\begin{warningbox}{El riesgo del periodo de espera: no es equivocarse en el juicio, sino no resistir hasta que ocurra}
Los emprendedores con frecuencia alcanzan a ver una dirección, pero no el momento. El juicio sobre la dirección debe corresponder con el ciclo de efectivo, la paciencia del equipo y el grado de terminación del producto. Llegar demasiado pronto desgasta; llegar demasiado tarde hace perder la ventana.
\end{warningbox}

\subsection{Resistir la retroalimentación positiva ajena a la estrategia}

Esta subsección añade un detalle. Durante la espera, los desarrollos a la medida para grandes clientes pueden dar retroalimentación positiva al equipo; el equipo de ventas se entusiasma y el efectivo de corto plazo también resulta atractivo. Pero si los proyectos a la medida se alejan del producto SaaS estándar, absorben recursos de desarrollo y afectan la iteración de la línea principal. Lo que hizo Xiao Hong fue rechazar la retroalimentación positiva ajena a la estrategia y reservar los recursos para el producto que «pudiera atender la demanda cuando ocurriera el evento». Este juicio es más concreto que el «enfoque» habitual: no se trata de rechazar todos los ingresos, sino de rechazar los ingresos que destruirían la espera estratégica.

\begin{importantbox}{Experiencia práctica: no toda retroalimentación positiva debe aceptarse}
Si una retroalimentación proviene de una dirección no estratégica, puede ser más peligrosa que una retroalimentación negativa. La retroalimentación negativa te pone en alerta; la retroalimentación positiva ajena a la estrategia lleva al equipo a creer erróneamente que encontró un nuevo camino y termina desviándolo de la ventana que en realidad estaba esperando.
\end{importantbox}

\subsection{Resumen de la sección}

Anticipar lo que anticiparán las grandes empresas fue el método central del primer emprendimiento de Xiao Hong: entender cómo regulará una plataforma su ecosistema, preparar con anticipación una alternativa que cumpla con la normativa y resistir las tentaciones de corto plazo durante la espera. Este método se trasladó después a la IA: anticipar cuándo aparecerán las capacidades de los modelos y colocar de antemano ahí el contenedor de la aplicación.
\section{Monica: de la extensión de navegador a la plataforma de aplicaciones de IA}

Lo anterior trató el método del primer emprendimiento; esta sección pasa al segundo. A finales de 2022, el equipo de Xiao Hong (肖弘) vio el potencial de la IA generativa y empezó a construir Monica. Al principio eligió una extensión de navegador, y no una chatbot app independiente, porque la extensión se integra con más facilidad en el flujo de trabajo que el usuario ya tiene; después, la adquisición de ChatGPT for Google le dio a Monica tráfico de arranque en frío y retroalimentación de usuarios.

\begin{figure}[H]
\centering
\includegraphics[width=0.96\textwidth]{monica-product-evolution.png}
\caption{Evolución del producto Monica: de la extensión de navegador y ChatGPT for Google a la plataforma de aplicaciones de IA. Diagrama conceptual propio, elaborado a partir de la conversación de 01:05:55--01:48:38.}
\end{figure}

\begin{knowledgebox}{Lectura de la figura: la extensión no es el destino final, pero es la puerta de entrada al flujo de trabajo}
De Plugin a ChatGPT for Google, Monica.im, la retroalimentación de usuarios y el Agent, la figura destaca el arranque en frío y la integración en el flujo de trabajo. La extensión hace que la IA aparezca donde el usuario ya trabaja, en lugar de pedirle que abra por iniciativa propia otra ventana de chat.
\end{knowledgebox}

\subsection{Por qué no construir directamente una Chatbot App}

Esta subsección explica la elección de la forma del producto. Una chatbot app independiente tiene que responder a la pregunta «¿por qué el usuario te abre a ti y no a la ChatGPT original?»; la extensión, en cambio, puede aparecer en escenarios como la búsqueda, la lectura de páginas web, la escritura, la traducción y el resumen, lo que reduce la fricción de inicio. Xiao Hong también reconoce que, antes del lanzamiento de la App oficial de ChatGPT, hubo una ventana de oportunidad para las apps que aprovecharon algunos productos extranjeros que solo envolvían el modelo; sin embargo, Monica optó por entrar al flujo de trabajo a través de la extensión.

\begin{knowledgebox}{Términos clave: arranque en frío, flujo de trabajo, contenedor}
\noindent\begin{tabularx}{\textwidth}{p{0.18\textwidth}p{0.36\textwidth}X}
\toprule
Término & Significado & Función en este episodio \\
\midrule
Arranque en frío & Problema de arranque cuando el producto, en su etapa temprana, carece de usuarios y retroalimentación & ChatGPT for Google le aportó usuarios y retroalimentación a Monica. \\
Flujo de trabajo & Secuencia de pasos y entorno de herramientas con los que el usuario completa una tarea & La extensión puede integrarse en el flujo de trabajo existente. \\
Contenedor & Forma de producto que da cabida a las capacidades del modelo & Un buen contenedor convierte las capacidades del modelo en tareas utilizables. \\
Comprensión del usuario & Conocimiento profundo de distintas regiones, culturas y formas de uso & Un producto para mercados extranjeros no puede limitarse a traducir la interfaz. \\
\bottomrule
\end{tabularx}
\end{knowledgebox}

\subsection{Por qué rehacer Monica en lugar de cambiarle el nombre a ChatGPT for Google}

Esta subsección conserva una decisión de producto clave. ChatGPT for Google tenía crecimiento orgánico y búsquedas de marca, pero su nombre parecía más el de una función que el de una marca a largo plazo. Cambiarle el nombre directamente podía dañar el crecimiento; y si el equipo se limitaba a mantener el producto anterior, no aprendería crecimiento en el extranjero, compra de publicidad, KOL ni comprensión del usuario. Por eso el equipo rehízo Monica y conservó el producto anterior como colchón de seguridad. Así contaba con activos de arranque en frío y, a la vez, obtenía entrenamiento real en crecimiento.

\begin{importantbox}{Nota de clase: el crecimiento comprado no sustituye el aprendizaje del equipo}
Adquirir un producto con tráfico puede acortar el arranque en frío; pero si el equipo se apoya únicamente en el crecimiento orgánico, no aprende crecimiento ni comprensión del usuario. Rehacer Monica tenía el sentido de devolver al equipo a un entorno en el que tuviera que conseguir la retroalimentación por sí mismo.
\end{importantbox}

\subsection{Salir al extranjero: el valor antes que la técnica}

Esta subsección examina la elección del mercado extranjero. Xiao Hong considera que el mercado nacional es un gran mercado unificado: ganarlo significa mucho, pero la competencia es feroz; el extranjero se parece más a muchos territorios pequeños, con idiomas, monedas, husos horarios, estéticas y culturas distintas, lo que deja más espacio local a las empresas emergentes. Que Monica eligiera empezar por el extranjero requirió valor: el equipo no había vivido mucho tiempo fuera y el inglés no era su fortaleza principal, pero la ventana de oportunidad importaba más que una preparación perfecta.

\begin{warningbox}{Salir al extranjero no es traducir la interfaz}
Un producto para el extranjero debe considerar varios idiomas, varias monedas, husos horarios, pagos, cumplimiento normativo, estética y hábitos de los usuarios. La IA puede ayudar a traducir, pero no resuelve por sí sola la comprensión del usuario. Salir de verdad al extranjero significa que la organización aprende de forma continua sobre mercados distintos.
\end{warningbox}

\subsection{Resumen de la sección}

La trayectoria de Monica muestra que emprender con aplicaciones de IA no consiste en envolver un modelo sin más, sino en elegir la entrada adecuada, la forma de arranque en frío, el nombre de marca, el alcance del mercado y el mecanismo de retroalimentación. Monica trasladó la experiencia del primer emprendimiento en comercialización, crecimiento y anticipación de plataformas a la ventana de oportunidad de las aplicaciones de IA.
\section{Metodología de aplicaciones de IA: capacidad del modelo, complemento de escenarios y desbordamiento}

La sección anterior trató sobre Monica; esta sección extrae la clasificación de aplicaciones de IA de Xiao Hong (肖弘). Él toma las aplicaciones de IA que se han vuelto populares como un número reducido de datos y resume tres tipos de oportunidades: primero, complementar escenarios principales, por ejemplo mejorar la experiencia en escenarios ya existentes como la búsqueda, la escritura, el navegador o la ofimática; segundo, nuevas formas de uso que trae la capacidad de los modelos, que tienden a despegar sobre todo en la generación de imágenes y video; tercero, el desbordamiento de la capacidad del modelo hacia un dominio específico, donde ya resuelve el problema central de ese dominio y solo falta un buen contenedor que la reciba, como Cursor, que recibe la capacidad de coding.

\begin{figure}[H]
\centering
\includegraphics[width=0.92\textwidth]{ai-app-taxonomy.png}
\caption{Método de clasificación de aplicaciones de IA: complemento de escenarios principales, cambio en la capacidad del modelo y desbordamiento hacia dominios específicos. Diagrama conceptual propio, elaborado a partir de la conversación de 02:02:14--02:09:32.}
\end{figure}

\begin{knowledgebox}{Lectura de la figura: buscar oportunidades no es perseguir modas, sino encontrar la brecha entre capacidad y escenario}
El complemento de escenarios principales observa dónde le falta IA al flujo de trabajo actual; el cambio de capacidad observa las nuevas formas de uso que abre un modelo nuevo; el desbordamiento hacia un dominio observa si un modelo general ya resolvió la parte más difícil de un dominio vertical, de modo que la empresa de producto solo necesita completar el flujo de trabajo periférico.
\end{knowledgebox}

\subsection{El ejemplo de Cursor: desbordamiento de la capacidad central y contenedor}

Esta subsección desarrolla el caso de Cursor. Xiao Hong cuenta que antes generaba código conversando con una IA y luego tenía que copiar y pegar a mano, crear archivos nuevos y comparar diferencias; después se dio cuenta de que lo que Cursor hizo bien fue convertir en ingeniería esas acciones periféricas. La capacidad central de escribir código proviene del foundation model, y la empresa de producto no necesita resolverla de nuevo; pero sí puede ofrecer el contenedor del editor, las operaciones con archivos, la comparación de diferencias, la gestión del contexto y una interacción fluida. Este es el ejemplo típico de «una vez que la capacidad del modelo se desborda, un contenedor la recibe».

\begin{importantbox}{Fórmula de oportunidad para aplicaciones de IA}
Si un modelo general ya resolvió la capacidad central de un dominio, pero el usuario todavía tiene que completar a mano muchos pasos periféricos, la empresa de aplicaciones puede construir el contenedor: conectar la capacidad del modelo con el flujo de trabajo real y eliminar fricciones como copiar, pegar, comparar, invocar herramientas y guardar resultados.
\end{importantbox}

\subsection{Fabricantes de modelos grandes frente a empresas de aplicaciones}

Esta subsección distingue dos tipos de empresas. Los fabricantes de modelos pueden decidir cuándo aparece una capacidad, por lo que tienen mayor control; las empresas de aplicaciones no pueden decidir los límites de la capacidad de los modelos y solo pueden anticipar, esperar e integrarla con rapidez. Pero las empresas de aplicaciones también tienen ventajas: pueden usar los modelos de todo el mundo y concentrarse en escenarios específicos, la experiencia de usuario, la marca, el crecimiento y la retroalimentación. Xiao Hong lo ve como la ruta de «comercio, industria, tecnología» (贸工技): primero construir usuarios y producto, y después completar gradualmente la ingeniería y las capacidades de modelo específicas.

\begin{figure}[H]
\centering
\includegraphics[width=0.92\textwidth]{model-company-vs-app-company.png}
\caption{Fabricantes de modelos grandes frente a empresas de aplicaciones: los fabricantes construyen capacidad general; las empresas de aplicaciones construyen escenarios, contenedores y retroalimentación. Diagrama conceptual propio, elaborado a partir de la conversación de 02:21:01--02:23:51.}
\end{figure}

\begin{knowledgebox}{Lectura de la figura: controlar cuándo surge la capacidad frente a recibir su desbordamiento}
Los fabricantes de modelos pueden impulsar la capacidad de base e incluso definir cuándo aparece una capacidad; las empresas de aplicaciones se parecen más a quien está al borde de la ola de capacidades y coloca de antemano el contenedor de producto. La relación entre ambas no es un simple vínculo de proveedor y cliente, sino distintas combinaciones de control, velocidad, comprensión del usuario y profundidad en el escenario.
\end{knowledgebox}

\subsection{Lo que enseña DeepSeek}

Esta subsección examina el impacto de DeepSeek en el ecosistema chino de aplicaciones de IA. Xiao Hong dice que DeepSeek fue «la más desapegada y, aun así, la que obtuvo los mejores resultados», lo que dio un impulso anímico a los emprendedores de aplicaciones: menos actuación y más capacidad real. El éxito repentino de DeepSeek también lo llevó a replantearse la difusión de productos: OpenAI o1 tiene una gran capacidad, pero su forma de presentarse como producto no logró captar la atención del gran público, mientras que DeepSeek llegó a usuarios de todo el mundo con una postura más modesta y una forma de uso más directa.

\begin{figure}[H]
\centering
\includegraphics[width=0.90\textwidth]{deepseek-product-lesson.png}
\caption{Lo que enseña DeepSeek como producto: el equipo más desapegado obtuvo los mejores resultados y dio un impulso anímico a los emprendedores de aplicaciones. Diagrama conceptual propio, elaborado a partir de la conversación de 02:23:51--02:29:00.}
\end{figure}

\begin{knowledgebox}{Lectura de la figura: la capacidad y la postura cambian juntas la difusión}
Del lado de DeepSeek se destacan la capacidad, el código abierto y la postura modesta; del lado del emprendimiento de aplicaciones se destaca el valor real. La lección aquí no es que todas las empresas deban adoptar una actitud desapegada, sino que la difusión de un producto no puede depender solo de una postura de marketing: la capacidad de fondo y la percepción del usuario deben sostenerse por sí mismas.
\end{knowledgebox}

\subsection{DeepSeek frente a o1: además de tener gran capacidad, el usuario tiene que percibirla}

Esta subsección añade la perspectiva de producto de la segunda entrevista. Cuando Xiao Hong habla de DeepSeek, no está evaluando un modelo, sino analizando cómo un acontecimiento de capacidad de un modelo se convierte en un acontecimiento en la mente de los usuarios. OpenAI o1 logró un avance importante en capacidad de razonamiento, pero su presentación como producto se inclinaba más hacia «una versión del modelo» y «capacidad especializada»; DeepSeek, en cambio, entró en el campo de visión del usuario común con una barrera de entrada más baja, más apertura y mayor facilidad de difusión. Para los emprendedores de aplicaciones, esto les recuerda que la capacidad técnica por sí sola no equivale automáticamente al éxito de un producto: la capacidad tiene que presentarse como una experiencia que el usuario pueda entender, probar y difundir directamente.

\begin{importantbox}{Los tres niveles de la presentación de un producto}
El primer nivel es que la capacidad exista realmente; el segundo, que el usuario pueda acceder a ella y verificarla a bajo costo; el tercero, que el usuario pueda explicársela a otros en una frase. Si DeepSeek dio un impulso anímico a los emprendedores de aplicaciones, es porque logró articular al mismo tiempo la capacidad, la postura abierta y la vía de difusión.
\end{importantbox}

\noindent\begin{tabularx}{\textwidth}{p{0.20\textwidth}p{0.34\textwidth}X}
\toprule
Dimensión & Presentación al estilo de o1 & Presentación al estilo de DeepSeek \\
\midrule
Percepción del usuario & Se parece más a una versión de modelo de razonamiento avanzado & Se parece más a un acontecimiento de capacidad que cualquiera puede probar. \\
Forma de difusión & Mayor presencia en círculos especializados y discusiones de benchmark & Mayor difusión por el código abierto, el bajo costo y la comunidad en chino. \\
Lección para las aplicaciones & Gran capacidad, pero la presentación como producto no siempre es suficiente & La capacidad combinada con la distribución da confianza a la capa de aplicaciones. \\
\bottomrule
\end{tabularx}

\subsection{Resumen de la sección}

La metodología de aplicaciones de IA no consiste en «tomar un modelo y ponerle una capa encima». Exige juzgar el siguiente paso de la capacidad de los modelos, encontrar el escenario principal del usuario, eliminar las fricciones del contenedor, comprender al usuario de forma continua y aceptar que la forma del producto se reescribirá una y otra vez conforme avance la capacidad de los modelos.
\section{Manus: el primer disparo de los Agent en China}

Las secciones anteriores trataron sobre Monica y la metodología de aplicaciones; esta sección entra en Manus. La segunda entrevista ocurrió antes del lanzamiento de Manus, y Xiao Hong (肖弘) lo describió como un producto nuevo a punto de lanzarse, el primer disparo de los Agent en China. Su entusiasmo provenía de una especie de A-ha moment: ya no parecía que estuviera haciendo una herramienta común, sino que estaba creando una especie de vida capaz de actuar. Lo importante aquí no es divinizar al Agent, sino entender que la forma del producto pasa de «responder preguntas» a «ejecutar tareas».

\begin{figure}[H]
\centering
\includegraphics[width=0.96\textwidth]{manus-agent-thesis.png}
\caption{Manus: el primer disparo de los Agent en China, de la tarea, las herramientas, el entorno y la memoria hasta un resultado entregable. Diagrama conceptual de elaboración propia, basado en la conversación de 02:31:12--02:49:27.}
\end{figure}

\begin{knowledgebox}{Lectura de la figura: el ciclo mínimo de un producto Agent}
De Task a Tools, Environment, Memory y Deliverable: un Agent no es solo una ventana de chat. El usuario da un objetivo, el sistema descompone la tarea, invoca herramientas, actúa en el entorno, conserva el estado y, al final, entrega un resultado. Si falta cualquiera de los eslabones, se vuelve a un asistente común o a un script de flujo de trabajo.
\end{knowledgebox}

\subsection{Por qué surge la sensación de «crear vida»}

Esta subsección explica el A-ha moment. Una herramienta común ejecuta una función de un solo paso según la instrucción del usuario; un Agent, en torno a un objetivo, descompone la tarea por iniciativa propia, la ejecuta en varios pasos, lee la retroalimentación y se corrige a sí mismo. Lo que el usuario ve no es una respuesta, sino un sistema que hace avanzar las cosas en su lugar. Esa acción continua, que supera las expectativas, hace que el producto pase de sentirse como una herramienta a sentirse como algo vivo. Pero al mismo tiempo exige más controlabilidad, más explicabilidad y capacidad de recuperarse de fallas.

\begin{figure}[H]
\centering
\includegraphics[width=0.86\textwidth]{aha-lifeform.png}
\caption{A-ha Moment: el producto Agent pasa de sentirse como una herramienta a sentirse como algo vivo que actúa. Diagrama conceptual de elaboración propia, basado en la conversación de 02:47:49--02:49:27.}
\end{figure}

\begin{warningbox}{Sensación de vida no significa renunciar al control}
Cuanto más parece un Agent «capaz de hacer cosas», más necesita permisos, confirmaciones, auditoría, posibilidad de interrupción y recuperación de errores. Que el producto sorprenda es solo el primer paso; el uso a largo plazo exige que sea controlable y confiable.
\end{warningbox}

\subsection{El mecanismo de Manus: tarea, entorno y entregable}

Esta subsección desglosa la lógica de producto de Manus. Un asistente de IA común se queda sobre todo en la conversación y las sugerencias; Manus busca que el usuario le encomiende una tarea, que entre en un entorno operable, invoque el navegador, archivos, búsqueda, código o herramientas externas, avance de forma continua y, al final, entregue un resultado. Esto significa que el centro del diseño del producto se desplaza de la «experiencia de conversación» al «ciclo de vida de la tarea»: cómo se entiende la tarea, cómo se hace visible el plan, cómo se controlan las invocaciones de herramientas, cómo se guardan los estados intermedios, cómo se recupera de las fallas y cómo el usuario puede dar por aceptado el entregable.

\begin{knowledgebox}{Términos clave: el ciclo de vida de la tarea en un producto Agent}
\noindent\begin{tabularx}{\textwidth}{p{0.20\textwidth}p{0.34\textwidth}X}
\toprule
Etapa & Función & Riesgo de producto \\
\midrule
Comprensión de la tarea & Convertir el objetivo del usuario en un plan ejecutable & Malinterpretar el objetivo hace que todo lo posterior salga mal. \\
Invocación de herramientas & Operar páginas web, archivos, código y API & Los permisos, el costo y los límites de seguridad deben estar claros. \\
Estado intermedio & Guardar el plan, la evidencia, las fallas y los productos & Perder el estado rompe las tareas largas. \\
Entregable & Convertir el resultado de las acciones en una salida aceptable & El usuario debe poder revisar y modificar el resultado. \\
\bottomrule
\end{tabularx}
\end{knowledgebox}

\begin{warningbox}{Un Agent no es «un Chatbot más largo»}
Si solo da respuestas más largas y maneja más contexto, sigue siendo un asistente. La línea divisoria del Agent está en si puede entrar en un entorno a actuar y consolidar el resultado de esas acciones en un entregable que el usuario pueda aceptar.
\end{warningbox}

\subsection{Cuando las grandes empresas te entienden, llega el momento de peligro}

Esta subsección examina la competencia. Xiao Hong dice que el momento en que las grandes empresas pueden entender tu innovación es un momento muy peligroso. Esto se debe a que la ventaja de las empresas de aplicaciones en etapa temprana suele provenir de los huecos que las grandes empresas todavía no han notado, no pueden cubrir rápidamente o, por estrategia, no quieren cubrir. Una vez que una gran empresa lo entiende y decide entrar, la empresa de aplicaciones ya debe haber construido presencia en la mente del usuario, flujos de trabajo, datos, marca o velocidad organizacional; de lo contrario, la plataforma y las capacidades del modelo la van a comprimir.

\begin{importantbox}{La ventana de oportunidad para emprender con Agent}
El mejor momento para una empresa de aplicaciones es cuando la capacidad del modelo ya es suficiente, la necesidad del usuario empieza a manifestarse, pero las grandes empresas todavía no entienden por completo la forma del producto. La ventana no dura mucho tiempo; hay que consolidar barreras de entrada antes de que las grandes empresas lo entiendan.
\end{importantbox}

\subsection{Lo que las grandes empresas no hacen, lo hacen las empresas de aplicaciones}

Esta subsección conecta Monica con Manus. Xiao Hong mencionó varias veces que, por razones estratégicas y organizacionales, las grandes empresas no hacen ciertas cosas: una empresa de modelos no necesariamente quiere integrar todos los modelos, una empresa de plataforma no necesariamente quiere hacerse cargo de todos los flujos de trabajo marginales, y el chatbot del propio fabricante no necesariamente va a profundizar en cada escenario vertical. La oportunidad de las empresas de aplicaciones suele estar justamente en esos huecos donde «el usuario tiene la necesidad, las grandes empresas no lo harán a corto plazo y la capacidad del modelo ya es suficiente». La hipótesis de Manus es que la ejecución de tareas de propósito general ya entró en la ventana en la que puede convertirse en producto, mientras que las grandes empresas todavía no han convertido por completo esa experiencia en un producto masivo.

\begin{importantbox}{Cómo juzgar la ventana para una empresa de aplicaciones}
Para saber si vale la pena aprovechar una ventana, se pueden plantear tres preguntas: si el usuario ya tiene un problema claro; si el modelo ya resolvió la capacidad central más difícil; si las grandes empresas, por razones de estrategia, ritmo u organización, por ahora no van a profundizar en ello. Si las tres se cumplen a la vez, la empresa de aplicaciones tiene espacio para llegar primero.
\end{importantbox}

\subsection{Resumen de la sección}

La importancia de Manus radica en que lleva el emprendimiento de aplicaciones de IA en China un paso más allá, de las herramientas, los plugins y los asistentes hacia los Agent. Lo que pone a prueba no es solo la capacidad del modelo, sino si el producto puede gestionar tareas, herramientas, entorno, memoria, permisos y resultados entregables.
\section{Mentalidad de Founder: no linealidad, teoría de juegos y la edad de la época}

Las secciones anteriores trataron sobre el producto y los Agent; esta sección concluye con la mentalidad de Founder. Xiao Hong (肖弘) insiste en que el mundo no evoluciona por extrapolación lineal. El emprendedor no puede limitarse a razonar lógicamente con las condiciones actuales; debe pensar en términos de juego estratégico: cambiar la situación, influir en las expectativas de los demás, elegir el momento y convertirse en una variable relevante. Esta idea atraviesa sus juicios sobre el capital, las plataformas, las capacidades de los modelos, DeepSeek, Manus y su propia etapa de vida.

\begin{figure}[H]
\centering
\includegraphics[width=0.90\textwidth]{founder-game-thinking.png}
\caption{El pensamiento estratégico del Founder: no extrapolar linealmente, sino convertirse en una variable relevante. Diagrama conceptual propio, elaborado a partir de la conversación entre 03:09:08 y 03:15:00.}
\end{figure}

\begin{knowledgebox}{Lectura de la figura: diferencia entre razonamiento lógico y pensamiento estratégico}
El razonamiento lógico busca la solución óptima con las condiciones dadas; el pensamiento estratégico cambia las condiciones mismas y obliga a los demás participantes a reevaluarte. El Founder no es un observador que pronostica, sino un actor dentro del juego.
\end{knowledgebox}

\subsection{Pensar con la edad de la época}

Esta subsección explica la «edad de la época». Xiao Hong nació en 1992 y vivió el final del internet móvil, el SaaS 2B, la IA generativa y la explosión de los Agent. Advierte que no hay que juzgar las oportunidades solo con la edad biológica, sino con la edad de la época: cuando empieza un nuevo ciclo tecnológico, todos pueden volver a la línea de salida. Quien emprende con aplicaciones de AI necesita identificar en qué etapa del ciclo tecnológico se encuentra: la de ventaja inicial, la de madurez o la de reestructuración.

\begin{importantbox}{Experiencia práctica: el emprendedor debe ubicarse en las coordenadas de su época}
Tener 32 años al final del internet móvil no es lo mismo que tenerlos en el primer año de los Agent. Para saber si uno llegó tarde, no basta con mirar la edad: hay que ver si las nuevas plataformas, los nuevos modelos y los nuevos comportamientos de los usuarios apenas empiezan a reescribir las reglas.
\end{importantbox}

\subsection{Por qué no desarrollar un modelo base}

Esta subsección responde a una pregunta importante. Xiao Hong no desarrolla un modelo base no por falta de respeto a las capacidades de los modelos, sino porque los recursos de cada empresa son distintos. Las empresas de modelos de primer nivel amplían la frontera de las capacidades y merecen respeto; las empresas de aplicaciones saben mejor cómo atender a los usuarios, aprovechar las capacidades de los modelos, construir flujos de trabajo y crear marca. A largo plazo, una empresa de aplicaciones puede entrenar modelos pequeños o modelos económicos para tareas específicas, pero no necesita apostar por una base general desde el day one.

\begin{knowledgebox}{Términos clave: fabricantes de modelos, empresas de aplicaciones, modelos económicos}
\noindent\begin{tabularx}{\textwidth}{p{0.20\textwidth}p{0.34\textwidth}X}
\toprule
Concepto & Significado & Papel en este episodio \\
\midrule
Fabricante de modelos & Empresa que entrena un foundation model de propósito general & Determina la frontera de capacidades de los modelos y cuándo aparecen. \\
Empresa de aplicaciones & Empresa que hace productos para escenarios concretos de usuario & Aprovecha el excedente de capacidades de los modelos y consolida escenarios y marca. \\
Modelo económico & Modelo pequeño, suficiente para una tarea específica y de costo controlable & La empresa de aplicaciones puede completar gradualmente capacidades de modelo puntuales. \\
Democratización tecnológica & Los modelos generales permiten a equipos pequeños lograr capacidades antes difíciles & Los desarrolladores independientes y los equipos pequeños obtienen nuevas oportunidades. \\
\bottomrule
\end{tabularx}
\end{knowledgebox}

\subsection{El Founder no tiene opción}

Esta subsección recoge el estado final del fundador. Xiao Hong dice que el Founder no tiene opción: una vez dentro del juego, no puede limitarse a tomar decisiones racionales como un espectador; muchas cosas hay que sostenerlas. La vida de emprendedor se vuelve cara: el estado psicológico, el costo de vida, la responsabilidad con el equipo, los compromisos con el capital y los cambios externos elevan el precio de cada decisión. El Founder tiene que enfrentar la codicia, la ira y la ignorancia, y también reconocer que la época, el capital y el deseo lo arrastran.

\begin{warningbox}{La mentalidad de Founder no es autoayuda motivacional}
«Convertirse en variable» no es una invitación a la impulsividad ciega, sino un recordatorio: el emprendedor debe atreverse a cambiar la situación y, al mismo tiempo, ver sus propios deseos, miedos y errores de juicio. El pensamiento estratégico sin autocrítica se convierte en mentalidad de apostador.
\end{warningbox}

\subsection{Llevar el pensamiento estratégico a la acción}

Esta subsección aterriza la mentalidad abstracta en acciones. El juego estratégico no es esoterismo; puede descomponerse en cuatro pasos: identificar a los participantes decisivos de la situación y evaluar sus motivaciones y restricciones; encontrar un cambio muy probable cuyo momento es incierto; prepararse con anticipación dentro de un costo tolerable; y, cuando el suceso ocurra, amplificarlo rápidamente con acciones de producto, difusión y organización. El SCRM para WeCom fue una jugada de este tipo, y Monica y Manus siguen una lógica similar: primero evaluar las capacidades de los modelos y los movimientos de las grandes empresas, y luego colocar con anticipación un contenedor de producto.

\noindent\begin{tabularx}{\textwidth}{p{0.18\textwidth}p{0.34\textwidth}X}
\toprule
Paso & Pregunta concreta & Ejemplo en EP95 \\
\midrule
Observar a los jugadores & Quién cambiará la situación y quién tiene poder de plataforma & Tencent combate los complementos no autorizados; los fabricantes de modelos liberan nuevas capacidades. \\
Observar las motivaciones & Por qué la otra parte tarde o temprano lo hará & La plataforma necesita regular su ecosistema; las empresas de modelos necesitan impulsar las capacidades. \\
Prepararse con anticipación & Con qué costo puedo tomar posición primero & Hacer primero el SCRM para WeCom; hacer primero los complementos de AI y Manus. \\
El suceso ocurre & Cómo amplificar rápidamente la retroalimentación positiva & Publicar un artículo la misma noche del bloqueo de los complementos no autorizados; acelerar los Agent después de DeepSeek. \\
\bottomrule
\end{tabularx}

\begin{knowledgebox}{Nota de clase: el pensamiento estratégico y la comprensión del usuario no se contraponen}
El pensamiento estratégico sirve para juzgar la ventana y la posición; la comprensión del usuario sirve para hacer bien el producto. Solo estrategia sin usuarios se vuelve especulación; solo usuarios sin estrategia puede hacer que se pierda una ventana estructural.
\end{knowledgebox}

\subsection{Resumen de la sección}

La mentalidad de Founder de Xiao Hong consiste en ver el emprendimiento como un juego dinámico: cambian las capacidades de los modelos externos, cambian los juicios de las grandes empresas, cambia la percepción de los usuarios y cambian las expectativas del capital. La tarea del Founder no es predecir el futuro de forma lineal, sino ganar peso como variable dentro de un sistema inestable.
\section{Síntesis y ampliación}

Esta sección condensa todo el episodio en un marco reutilizable. El hilo central de EP95 es cómo actúa un emprendedor de aplicaciones de IA en un ciclo tecnológico inestable: primero aprende de emprendimientos anteriores sobre comercialización, capital y anticipación de las plataformas; después aprovecha con Monica la oportunidad de los flujos de trabajo de la era de ChatGPT; luego apuesta con Manus por la ejecución de tareas en la era de los Agents; por último, usa el pensamiento de teoría de juegos para entender el papel del Founder en un mundo no lineal.

\begin{importantbox}{Siete conclusiones centrales}
Primera: una empresa de aplicaciones de IA no puede limitarse a seguir las tendencias de los modelos; debe llevar la capacidad del modelo a escenarios reales. Segunda: los productos de herramientas deben comercializarse pronto; la economía unitaria es la condición previa de un crecimiento sano. Tercera: el VC es una herramienta costosa; el Founder no puede tratar a los inversionistas como sus jefes. Cuarta: anticipar lo que anticipan las grandes empresas es el método decisivo para que una empresa de aplicaciones encuentre su ventana de oportunidad. Quinta: la capacidad de los modelos tenderá a igualarse tecnológicamente, pero el contenedor, la marca, la comprensión del usuario y la retroalimentación siguen siendo barreras competitivas. Sexta: Manus representa el salto de producto del asistente de IA al Agent. Séptima: el mundo no se extrapola de forma lineal; el Founder debe convertirse en una variable importante dentro del juego estratégico.
\end{importantbox}

\subsection{Tabla de consulta rápida de términos de este episodio}

\noindent\begin{tabularx}{\textwidth}{p{0.22\textwidth}p{0.34\textwidth}X}
\toprule
Término/producto & Significado en este episodio & Aspecto por observar \\
\midrule
Monica & all-in-one AI assistant, que pasó de extensión de navegador a puntos de entrada en múltiples plataformas & Integración en flujos de trabajo, crecimiento en el extranjero, comprensión del usuario. \\
ChatGPT for Google & Producto paralelo adquirido que aportó arranque en frío y retroalimentación de usuarios & Cómo un activo comprado ayuda a lanzar un producto nuevo. \\
Manus & Descrito antes de su lanzamiento como el primer disparo de los Agents en China & Capacidad de tareas, herramientas, entorno, memoria y entrega. \\
DeepSeek & Acontecimiento de modelos que, después del Año Nuevo chino, reconfiguró el ecosistema de base de las aplicaciones de IA en China & Capacidad, código abierto, perfil bajo y difusión del producto. \\
Cursor & Ejemplo representativo de cómo la capacidad del modelo se extiende al escenario de coding & Cómo el contenedor aprovecha la capacidad del modelo base. \\
VC & Herramienta de capital de alto costo & Usar el capital sin dejarse conducir por él. \\
Juego estratégico del Founder & Cambiar las condiciones, influir en las expectativas, convertirse en variable & Marco de acción en un mundo no lineal. \\
\bottomrule
\end{tabularx}

\subsection{Preguntas para el seguimiento}

Por último, esta sección convierte la entrevista en una lista de seguimiento. El valor de EP95 no está solo en registrar la víspera del lanzamiento de Manus, sino en ofrecer un marco de preguntas para evaluar a las empresas de aplicaciones de IA: si la capacidad de los modelos seguirá extendiéndose, si la empresa de aplicaciones tiene de verdad capacidad de contenedor y si el Founder puede consolidar barreras competitivas antes de que entren las grandes empresas.

\begin{enumerate}
\item ¿Podrá Manus convertir el A-ha moment previo al lanzamiento en flujos de trabajo de Agent estables, con alta retención y controlables?
\item ¿Podrá la ruta all-in-one assistant de Monica mantener su diferenciación frente a las aplicaciones de las grandes empresas y los Agents verticales?
\item ¿La barrera competitiva de las empresas de aplicaciones de IA terminará apoyándose más en la marca, los datos, los flujos de trabajo, la memoria del Agent o en modelos pequeños para tareas específicas?
\item Después de DeepSeek, ¿se atreverán los emprendedores chinos de aplicaciones de IA a crear productos nativos, en lugar de solo esperar la capacidad de los modelos extranjeros?
\item Una vez que las grandes empresas comprendan el mercado y entren en él, ¿podrán las empresas de aplicaciones mantener la ventaja gracias a una comprensión más rápida del usuario y a un ciclo cerrado más sólido en su escenario?
\item ¿Cómo se equilibra el pensamiento de juego estratégico del Founder con la construcción de la organización a largo plazo, el cumplimiento normativo, la seguridad y la confianza de los usuarios?
\end{enumerate}

\subsection{Lecturas adicionales}

\begin{itemize}
\item Si le interesa la organización posterior de Manus y su situación antes de la adquisición, puede contrastarlo con EP128, la última entrevista a Peak/Manus.
\item Si le interesan los informes técnicos de Agents y los paradigmas de entrenamiento, puede contrastarlo con EP110, síntesis de informes técnicos de Agents.
\item Si le interesan los métodos de producto para aplicaciones de IA, puede contrastarlo con EP101 YouWare, Lovart special y EP97, informe trimestral de modelos grandes.
\item Si le interesan la extensión de la capacidad de los modelos y los límites de la capacidad multimodal, puede contrastarlo con EP102, entrevista sobre multimodalidad a Zhang Xiangyu (张祥雨).
\end{itemize}

\end{document}
